% MT064 / explanation — canonical editable content.
\textbf{(a) 联系是组合多个学习器，区别是训练怎样组织。} Bagging对训练记录重采样，各基模型分别训练，最后投票或平均。Boosting根据当前组合没处理好的部分训练下一轮，例如提高难分样本权重，再把新模型加入组合。因此bagging各模型训练可独立进行，boosting训练具有前后依赖。


这个区别不延伸成“boosting预测也只能串行”。模型一旦训练好，各弱学习器对新输入的输出仍可分别计算，最后加权汇总。两者都不保证永远比另一种更快或更准。


\textbf{(b) 按分段公式画Huber图。} 横轴是\(r=y-f(x)\)。在\(-1<r<1\)画\(L=r^2/2\)的抛物线，经过\((0,0)\)；在两端改画直线\(L=|r|-1/2\)。接点为\((-1,0.5)\)、\((1,0.5)\)，到\(\pm2\)时高度为1.5，两段斜率在接点相接。


\StudyImage{CS5489/exam-review/Midterm/assets/answer-huber.png}


它对大的输出异常值比平方损失稳健。半平方损失的斜率是\(r\)，误差越大，对拟合的拉动越强；Huber两端斜率大小固定为1。补充例子：\(r=10\)时，半平方损失的斜率为10，Huber只有1。


“损失没有上界，所以不稳健”这个判断不对。Huber的损失值仍会增长，但增长斜率受限；这里控制的是极端残差的影响，不要求损失有上界。完整分段导数及输入异常的条件见\StudyRef{Answers}{mt063}{MT063}。


\textbf{出处：} 2025A期中 Q7，原卷第4页（10分）


\textbf{答案依据：} 2025A期中 Q7 印刷题干；本题解为我们的推导，不采用扫描件的学生作答。
